% © 2026 Ing. David Jaroš — CC BY-NC-ND 4.0
%
% This work is licensed under a Creative Commons Attribution-NonCommercial-NoDerivatives
% 4.0 International License (CC BY-NC-ND 4.0).
%
% License History: Earlier drafts (up to v0.3) were released under CC BY 4.0.
% From v0.4 onward, all material is released under CC BY-NC-ND 4.0 to protect
% the integrity of the theoretical work during ongoing academic development.
%
% See LICENSE.md for full license text.

\ifdefined\INCLUDEMODE
\else
  \documentclass[12pt,a4paper]{article}
  \usepackage[a4paper, margin=2.5cm]{geometry}
  \usepackage{amsmath, amssymb, amsthm}
  \usepackage{mathtools}
  \usepackage{hyperref}
  \usepackage{xcolor}
  \usepackage{mdframed}
  \usepackage{booktabs}

  \newtheorem{theorem}{Theorem}[section]
  \newtheorem{lemma}[theorem]{Lemma}
  \newtheorem{proposition}[theorem]{Proposition}
  \newtheorem{corollary}[theorem]{Corollary}
  \theoremstyle{definition}
  \newtheorem{definition}[theorem]{Definition}
  \theoremstyle{remark}
  \newtheorem{remark}[theorem]{Remark}

  \newmdenv[backgroundcolor=yellow!20, linecolor=orange, linewidth=1.5pt]{gapbox}
  \newmdenv[backgroundcolor=green!15, linecolor=green!60!black, linewidth=1.5pt]{resultbox}
  \newmdenv[backgroundcolor=blue!8, linecolor=blue!50, linewidth=1.5pt]{insightbox}
  \newmdenv[backgroundcolor=red!8, linecolor=red!60, linewidth=1.5pt]{warnbox}

  \title{\textbf{Discrete Symmetry Formalization, Step 3:\\
                 Symmetry Breaking Catalogue ---\\
                 Physical vs.\ Effective/Dissipative}}
  \author{David Jaro\v{s}}
  \date{2026-04-25}

  \begin{document}
  \maketitle

  \begin{abstract}
  We provide a complete catalogue of discrete symmetry breaking in the UBT
  framework, distinguishing two classes:
  (i) \emph{physical symmetry breaking} --- present in the fundamental action
  $S[\Theta]$ and with observable consequences;
  (ii) \emph{effective/dissipative symmetry breaking} --- arising in reduced
  descriptions of $\Theta$ relevant to AI/cryptographic models,
  not present as a fundamental violation.
  For each breaking type we identify the responsible term, state whether it is
  spontaneous or explicit, and state the UBT status (proved, candidate,
  gap, open).
  Open problems in the UBT discrete-symmetry sector are listed.
  \end{abstract}
\fi

%──────────────────────────────────────────────────────────────────────────────
\section{Classification of Symmetry Breaking}
\label{sec:cat:classification}
%──────────────────────────────────────────────────────────────────────────────

Symmetry breaking in the UBT context falls into three structurally distinct
categories:

\begin{description}
  \item[Explicit physical breaking]
    A term in $S[\Theta]$ does not respect the symmetry.  The breaking is
    present at the level of the fundamental action and has direct observable
    consequences (e.g., parity violation in weak interactions, CP violation
    in kaon decays).
  \item[Spontaneous physical breaking]
    The action is symmetric, but the ground state (vacuum solution $\Theta_0$)
    is not.  The symmetry is broken by the choice of vacuum, producing Goldstone
    bosons or massive gauge bosons (Higgs mechanism).
  \item[Effective/dissipative breaking]
    The fundamental action is symmetric, but a reduced or coarse-grained
    description of $\Theta$ (obtained by integrating out degrees of freedom or
    restricting to a particular solution sector) breaks the symmetry.
    This type of breaking is relevant for AI and cryptographic applications
    of UBT-inspired field equations, where dissipation and entropy production
    are modelled as field equations, not as fundamental violations.
\end{description}

\begin{warnbox}
\textbf{Critical distinction.}
Effective/dissipative symmetry breaking must \emph{never} be mistaken for a
prediction of fundamental $CPT$ violation.
The UBT full action satisfies $CPT$ invariance (as verified in Step~2).
Any apparent $CPT$ breaking in a specific model or application is an artefact
of the reduced description, not a prediction of the theory.
\end{warnbox}

%──────────────────────────────────────────────────────────────────────────────
\section{Physical Symmetry Breaking in UBT}
\label{sec:cat:physical}
%──────────────────────────────────────────────────────────────────────────────

\subsection{Parity Violation: $SU(2)_L$ Chiral Coupling}
\label{subsec:cat:parity}

\begin{definition}[Physical P breaking in UBT]
\label{def:cat:P_breaking}
Parity is explicitly broken in the UBT kinetic sector by the chiral structure
of the $SU(2)_L$ gauge coupling.  The relevant interaction term is:
\begin{equation}
  \mathcal{L}_{W} = \mathrm{Tr}\!\bigl[
    (P_L\Theta)^\dagger \,ig\tau^i W^i_\mu\, \partial^\mu(P_L\Theta)
  \bigr] + \mathrm{h.c.},
  \label{eq:cat:L_W}
\end{equation}
where $P_L = \tfrac{1}{2}(1 + P_2)$ projects onto the left-handed
(vector/odd) sector of $\mathcal{B}$ and $W^i_\mu$ are the $SU(2)_L$ gauge
bosons.
\end{definition}

\begin{proposition}[P-odd character of $\mathcal{L}_W$]
Under $P$ (Definition~2 of Step~1), $P_L \leftrightarrow P_R$, so:
\begin{equation}
  P[\mathcal{L}_W] = \mathrm{Tr}\!\bigl[
    (P_R\Theta)^\dagger \,ig\tau^i W^i_\mu\, \partial^\mu(P_R\Theta)
  \bigr] + \mathrm{h.c.} \neq \mathcal{L}_W.
\end{equation}
$\mathcal{L}_W$ is P-violating.  The violation is maximal in the sense that
there is no corresponding right-handed $W_R$ coupling in the UBT action at the
same energy scale.
\end{proposition}

\paragraph{UBT derivation status.}
The absence of $W_R$ coupling and the left-handed restriction of $W^i$ is
proved in \texttt{chirality/step3\_gap\_C1\_resolution.tex}: the
$W^\pm$ vertex in $S[\Theta]$ is $P_\psi$-odd (i.e., odd under $\psi \to -\psi$),
which implies $P$-odd coupling in the four-dimensional spacetime sense.
\textbf{Status: PROVED (conditional on algebraic completeness of the
$SU(2)_L$ derivation).}

\subsection{CP Violation: Complex Yukawa Sector}
\label{subsec:cat:CP}

\begin{definition}[CP violation in UBT]
\label{def:cat:CP_breaking}
CP violation arises from complex phases in the Yukawa coupling matrix $Y$
in the potential:
\begin{equation}
  V_Y = \mathrm{Tr}[\bar\Theta_L Y \Theta_R] + \mathrm{h.c.},
  \label{eq:cat:V_Yukawa}
\end{equation}
where $\Theta_L = P_L\Theta$, $\Theta_R = P_R\Theta$, and
$Y \in M_n(\mathbb{C})$ is the complex Yukawa matrix.
CP violation occurs when $\mathrm{Im}(\det Y) \neq 0$, i.e., when the
phases of $Y$ cannot be absorbed by field redefinitions.
\end{definition}

\begin{remark}[CKM analogue in UBT]
In the Standard Model, CP violation in the quark sector is described by the
CKM matrix with a single non-removable complex phase $\delta$.
In UBT, the analogous structure is the complex phase content of $Y$ acting
on the three generations encoded in the $\psi$-mode expansion
$\Theta = \sum_{n=1,2,3} \Theta_n e^{in\psi/R_\psi}$.
A complete UBT derivation of the CKM matrix from the $\psi$-sector is an
open problem (Gap~S3).
\end{remark}

\begin{gapbox}
\textbf{Gap S3 --- CKM phases from biquaternion $\psi$-modes:}
Express the Yukawa matrix $Y$ as a function of the complex time phases
$e^{in\psi/R_\psi}$ ($n = 1,2,3$) and derive the CKM mixing angles and
CP-violating phase from the $\Theta$-field dynamics.
This would constitute a \emph{derivation} of CP violation from first
principles within UBT, rather than a parametric input.
\textbf{Priority: HIGH.}
\end{gapbox}

\subsection{T Violation: $\theta$-Term and Strong CP Problem}
\label{subsec:cat:T}

\begin{definition}[Topological T breaking]
\label{def:cat:T_breaking}
The topological $\theta$-term (if present in the UBT action) breaks $T$:
\begin{equation}
  \mathcal{L}_\theta = \frac{\theta}{16\pi^2}\, G^a_{\mu\nu}\tilde{G}^{a\,\mu\nu},
  \label{eq:cat:L_theta}
\end{equation}
where $\tilde{G}^{\mu\nu} = \tfrac{1}{2}\epsilon^{\mu\nu\rho\sigma}G_{\rho\sigma}$
is the dual of the gluon field strength.
Under $T_{\mathrm{UBT}}$: $G\tilde{G} \to -G\tilde{G}$, so
$\mathcal{L}_\theta \to -\mathcal{L}_\theta$: explicit $T$ (and $P$) breaking.
This is the \emph{strong CP problem}: why is $\theta$ observed to be
$\lesssim 10^{-10}$ experimentally?
\end{definition}

\begin{gapbox}
\textbf{Gap S4 --- Strong CP problem in UBT:}
Does UBT predict $\theta = 0$ (or $\theta \ll 1$) from first principles, or
does it require an axion-like mechanism?
A natural candidate is that the $\psi$-integration in the biquaternionic
action dynamically averages $\theta$ to zero.
\textbf{Priority: MEDIUM.}
\end{gapbox}

\subsection{CPT Invariance: Theorem and Status}
\label{subsec:cat:CPT}

\begin{theorem}[UBT CPT Invariance]
\label{thm:cat:CPT}
The UBT action $S[\Theta]$ in the real sector ($\psi \to 0$ limit) is invariant
under $CPT = C \circ P \circ T_{\mathrm{UBT}}$:
\begin{equation}
  S[CPT[\Theta]] = S[\Theta].
\end{equation}
\end{theorem}

\begin{proof}[Proof sketch]
From Step~2, Table~\ref{tab:act:symmetries}: every term in
$\mathcal{L}_{\mathrm{UBT}}$ (including the $P$-violating chiral term and the
CP-violating complex Yukawa term) is $CPT$-invariant.
This is consistent with the Wightman--CPT theorem: any local, Lorentz-invariant
quantum field theory with a Hermitian Hamiltonian is $CPT$-invariant.
UBT satisfies all these conditions in the real sector.
\textbf{Status: Argued from general theorem + term-by-term verification
(Table~\ref{tab:act:symmetries}). Full biquaternionic sector: open (Gap~S2).}
\end{proof}

%──────────────────────────────────────────────────────────────────────────────
\section{Spontaneous Symmetry Breaking in UBT}
\label{sec:cat:spontaneous}
%──────────────────────────────────────────────────────────────────────────────

\subsection{Higgs Mechanism from $\Theta$-Vacuum}
\label{subsec:cat:higgs}

\begin{definition}[Spontaneous symmetry breaking via $\Theta_0$]
\label{def:cat:Higgs}
Let $\Theta_0 \in \mathcal{B}$ be the vacuum (ground state) solution satisfying
$\delta S[\Theta_0] = 0$ with minimum energy.
A symmetry $G$ of $S[\Theta]$ is \emph{spontaneously broken} if
$G \cdot \Theta_0 \neq \Theta_0$.

In UBT, the electroweak symmetry $SU(2)_L \times U(1)_Y$ is broken to
$U(1)_{\mathrm{em}}$ by a vacuum expectation value:
\begin{equation}
  \langle\Theta_0\rangle = v \cdot e_0 + \delta\Theta,
  \quad v \approx 174\ \mathrm{GeV},
  \label{eq:cat:vev}
\end{equation}
where $e_0 \in \{1\} \subset \mathcal{B}$ is the scalar (parity-even) direction.
The fluctuation $\delta\Theta$ around $\Theta_0$ gives the massive $W$, $Z$
bosons and the Higgs boson.
\end{definition}

\begin{remark}[P and C status after spontaneous breaking]
After $SU(2)_L \times U(1)_Y \to U(1)_{\mathrm{em}}$:
\begin{itemize}
  \item $P$ (spatial parity) is already explicitly broken by the
    $SU(2)_L$ coupling (Section~\ref{subsec:cat:parity}).
  \item $C$ is broken by the charged-current structure.
  \item $CP$: broken by the CKM phase (Section~\ref{subsec:cat:CP}).
  \item $CPT$: preserved (Theorem~\ref{thm:cat:CPT}).
\end{itemize}
\end{remark}

%──────────────────────────────────────────────────────────────────────────────
\section{Effective and Dissipative Symmetry Breaking}
\label{sec:cat:effective}
%──────────────────────────────────────────────────────────────────────────────

In AI-inspired and cryptographic applications of UBT-type field equations,
the field $\Theta$ is used as a model for complex adaptive systems.
The fundamental UBT action is replaced by effective equations that describe
irreversible dynamics, information flow, or computational state evolution.
These equations generically break the discrete symmetries $C$, $P$, $T$ at
the \emph{effective level}, without implying any violation of the
fundamental $CPT$ theorem.

\subsection{Explicit Time Asymmetry}
\label{subsec:cat:time_asym}

\begin{definition}[Dissipative time evolution]
\label{def:cat:dissipation}
An effective dissipative term takes the form:
\begin{equation}
  \mathcal{L}_{\mathrm{diss}}
  = -\kappa_T\,\mathrm{Re}\,\mathrm{Tr}\!\bigl[
      \Theta^\dagger (\partial_t \Theta)
    \bigr],
  \label{eq:cat:L_diss}
\end{equation}
with $\kappa_T > 0$ a phenomenological damping coefficient.
\end{definition}

\begin{proposition}
$\mathcal{L}_{\mathrm{diss}}$ breaks $T_{\mathrm{UBT}}$: under $t \to -t$,
$\partial_t \to -\partial_t$, so
$\mathcal{L}_{\mathrm{diss}} \to +\kappa_T \mathrm{Re}\,\mathrm{Tr}[\Theta^\dagger \partial_t\Theta]
= -\mathcal{L}_{\mathrm{diss}}$.
\end{proposition}

\paragraph{Physical interpretation.}
$\mathcal{L}_{\mathrm{diss}}$ does not appear in the fundamental UBT action.
It arises when modelling the evolution of the $\Theta$-field in an environment
(open system), where irreversibility results from integrating out environmental
degrees of freedom.  This is entirely consistent with $CPT$ invariance at the
fundamental level: the environment that was integrated out is the ``time-reversed
universe.''

\subsection{Non-Hermitian Mass Terms}
\label{subsec:cat:non_herm}

\begin{definition}[Non-Hermitian mass]
\label{def:cat:non_hermitian}
An effective non-Hermitian mass term:
\begin{equation}
  V_{\mathrm{eff}} = (m^2 + i\gamma)\,\mathrm{Tr}[\Theta^\dagger\Theta],
  \quad \gamma \in \mathbb{R},
  \label{eq:cat:V_eff}
\end{equation}
with $\gamma \neq 0$ describes amplification ($\gamma < 0$) or damping
($\gamma > 0$) of the field.
\end{definition}

\begin{proposition}
$V_{\mathrm{eff}}$ with $\gamma \neq 0$ breaks $C$ (since complex conjugation
maps $i\gamma \to -i\gamma$, changing the sign of the imaginary part) and $T$
(since $T$ is real-time reversal which in an antiunitary formulation would
also map $i \to -i$).
$CPT$ is broken if $\gamma \neq 0$ and the field equations are not derived
from a real action.
\end{proposition}

\begin{warnbox}
Non-Hermitian mass terms \eqref{eq:cat:V_eff} are \textbf{not} part of the
fundamental UBT action $S[\Theta]$.  They are \textbf{effective descriptions}
of open-system dynamics (e.g., $\mathcal{PT}$-symmetric quantum mechanics,
gain-loss coupled waveguides, neural network dynamics).
Their $CPT$-breaking character is an artefact of the effective description,
not a prediction of UBT.
\end{warnbox}

\subsection{Entropy/Information Flow Terms}
\label{subsec:cat:entropy}

In information-theoretic and AI applications, UBT-inspired field equations
may include entropy-gradient terms:
\begin{equation}
  \mathcal{L}_{\mathrm{ent}}
  = -\beta \int d\psi\;
    \mathrm{Tr}\!\bigl[\Theta^\dagger \ln(\Theta^\dagger\Theta)\,\Theta\bigr],
  \label{eq:cat:L_ent}
\end{equation}
which models free-energy minimisation in a statistical ensemble.

\begin{proposition}
$\mathcal{L}_{\mathrm{ent}}$ is:
\begin{itemize}
  \item $P$-invariant (scalar under spatial reflection).
  \item $T_{\mathrm{UBT}}$-\emph{breaking}: entropy increases along the
    arrow of time, so the effective dynamics is time-asymmetric.
  \item $CPT$-invariant \emph{at the fundamental level}, but
    effectively $T$-breaking in the thermodynamic limit.
\end{itemize}
\end{proposition}

\paragraph{Separation principle.}
The arrow of time encoded in \eqref{eq:cat:L_ent} is an \emph{emergent},
thermodynamic phenomenon: it reflects the specific macroscopic state
(low-entropy initial condition) rather than a fundamental $T$-violation.
This is the UBT analogue of the second law of thermodynamics.

%──────────────────────────────────────────────────────────────────────────────
\section{Comparison Table: Physical vs.\ Effective Breaking}
\label{sec:cat:comparison}
%──────────────────────────────────────────────────────────────────────────────

\begin{table}[h]
\centering
\caption{Comparison of physical and effective symmetry breaking in UBT.}
\label{tab:cat:comparison}
\renewcommand{\arraystretch}{1.4}
\begin{tabular}{@{}p{4cm}p{3cm}p{3cm}p{4.5cm}@{}}
  \toprule
  Breaking & Symmetry broken & Type & UBT source \\
  \midrule
  $SU(2)_L$ chirality & $P$, $C$, $CP$ & Explicit physical &
    Chiral $W^\pm$ coupling in $\mathcal{L}_{\mathrm{kin}}$ \\
  CKM phases & $CP$, $C$ & Explicit physical &
    Complex Yukawa matrix $Y$ in $V(\Theta)$ \\
  Electroweak symmetry & $SU(2)_L \times U(1)_Y$ & Spontaneous physical &
    VEV $\langle\Theta_0\rangle = v\,e_0$ \\
  Strong CP ($\theta$-term) & $P$, $T$, $CP$ & Explicit physical (if $\theta \neq 0$) &
    $\mathcal{L}_\theta$ (not in minimal action; Gap S4) \\
  $\psi$-odd potential & $C$, $T$ (not $CPT$) & Explicit physical &
    $V_{\mathrm{odd}} = \mu\psi\,\mathrm{Tr}[\Theta^\dagger\Theta]$ \\
  \midrule
  Open-system dissipation & $T$ & Effective (open system) &
    $\mathcal{L}_{\mathrm{diss}} = \kappa_T \Theta^\dagger \partial_t\Theta$ \\
  Non-Hermitian mass & $C$, $T$ & Effective ($\mathcal{PT}$-sym.\  model) &
    $V_{\mathrm{eff}} = (m^2+i\gamma)\mathrm{Tr}[\Theta^\dagger\Theta]$ \\
  Entropy gradient & $T$ & Thermodynamic emergence &
    $\mathcal{L}_{\mathrm{ent}}$ (entropy minimisation) \\
  Gauge-choice asymmetry & depends & Gauge artefact &
    Non-invariant gauge fixing (not fundamental) \\
  \bottomrule
\end{tabular}
\end{table}

%──────────────────────────────────────────────────────────────────────────────
\section{Open Problems in the UBT Discrete Symmetry Sector}
\label{sec:cat:open}
%──────────────────────────────────────────────────────────────────────────────

The following open problems are identified in the discrete symmetry sector
of UBT, ordered by priority:

\begin{enumerate}
  \item \textbf{Gap S3 (HIGH)}: Derive the CKM matrix and CP-violating phase
    from the three-generation $\psi$-mode structure of $\Theta$.
    Expected location: intersection of
    \texttt{canonical/chirality/} and \texttt{canonical/interactions/}.
  \item \textbf{Gap S4 (MEDIUM)}: Determine whether UBT predicts
    $\theta_{\mathrm{QCD}} \approx 0$ from the biquaternionic action
    (potential resolution of the strong CP problem).
  \item \textbf{Gap S1 (MEDIUM)}: Resolve whether $T_{\mathrm{UBT}}$ should
    be antiunitary (composed with $P_1$) or unitary.
    Requires specifying the canonical sesquilinear form on $\Theta$-space.
  \item \textbf{Gap S2 (LOW)}: Verify $T_{\mathrm{UBT}}$ invariance of
    $\mathcal{L}_{\mathrm{grav}}$ in the full biquaternionic sector
    (including $\psi$-dependence of $\mathcal{G}_{\mu\nu}$).
  \item \textbf{Gap S5 (LOW)}: Classify all $\psi$-odd interaction terms
    that can appear in $V(\Theta)$ from symmetry arguments, and determine
    which are generated at loop level in the UBT effective field theory.
\end{enumerate}

\begin{resultbox}
\textbf{Summary of Step 3.}
\begin{itemize}
  \item \textbf{Physical $P$ breaking}: explicit, from chiral $SU(2)_L$ coupling
    (proved, see chirality derivation).
  \item \textbf{Physical $CP$ breaking}: explicit, from complex Yukawa matrix
    $Y$ (parametric; derivation from $\psi$-modes is open --- Gap S3).
  \item \textbf{Physical $T$ breaking}: potential, from $\theta$-term
    (not in minimal action; Gap S4).
  \item \textbf{$CPT$ invariance}: holds for all physical terms
    (consistent with CPT theorem).
  \item \textbf{Effective $T$ breaking}: arises in dissipative, non-Hermitian,
    or thermodynamic effective descriptions used in AI/crypto models.
    This is NOT a fundamental violation; it reflects the open-system or
    coarse-grained character of those models.
  \item \textbf{Key separation principle}: any apparent $CPT$ breaking in an
    effective model must be traced to an approximation (ignored environment,
    non-Hermitian coupling, truncated mode expansion) and does not falsify UBT.
\end{itemize}
\end{resultbox}

\ifdefined\INCLUDEMODE
\else
  \end{document}
\fi
